\documentclass[11pt,a4paper]{article}
\usepackage[utf8]{inputenc}
\usepackage[french]{babel}
\usepackage{amsmath}
\usepackage{amssymb}
\usepackage{geometry}
\usepackage{enumitem}
\usepackage{xcolor}
\usepackage{titlesec}
\usepackage{fancyhdr}
\usepackage{booktabs}
\usepackage{array}
\usepackage{listings}
\usepackage{tcolorbox}
\usepackage{algorithm}
\usepackage{algorithmic}

% Configuration de la page
\geometry{margin=2.5cm}
\pagestyle{fancy}
\fancyhf{}
\fancyhead[L]{\leftmark}
\fancyhead[R]{Exercices Machine Learning}
\fancyfoot[C]{\thepage}

% Configuration des titres
\titleformat{\section}
{\Large\bfseries\color{blue!70!black}}
{}
{0em}
{}[\titlerule]

\titleformat{\subsection}
{\large\bfseries\color{blue!50!black}}
{}
{0em}
{}

% Configuration pour le code Python
\lstset{
    language=Python,
    basicstyle=\ttfamily\small,
    keywordstyle=\color{blue}\bfseries,
    commentstyle=\color{gray}\itshape,
    stringstyle=\color{red},
    numbers=left,
    numberstyle=\tiny\color{gray},
    stepnumber=1,
    numbersep=5pt,
    backgroundcolor=\color{gray!10},
    frame=single,
    breaklines=true,
    showstringspaces=false,
    tabsize=4,
    morekeywords={np, pd, sklearn, plt, sns},
    inputencoding=utf8,
    extendedchars=true,
    literate={à}{{\`a}}1 {é}{{\'e}}1 {è}{{\`e}}1 {ê}{{\^e}}1 {ë}{{\"e}}1
             {ù}{{\`u}}1 {û}{{\^u}}1 {ü}{{\"u}}1 {ô}{{\^o}}1 {ç}{{\c c}}1
             {À}{{\`A}}1 {É}{{\'E}}1 {È}{{\`E}}1 {Ê}{{\^E}}1 {Ë}{{\"E}}1
             {Ù}{{\`U}}1 {Û}{{\^U}}1 {Ü}{{\"U}}1 {Ô}{{\^O}}1 {Ç}{{\c C}}1
}

% Configuration pour les zones importantes
\tcbuselibrary{breakable}
\newtcolorbox{questionbox}{
    colback=blue!5,
    colframe=blue!50,
    breakable,
    title=Question
}

\newtcolorbox{databox}{
    colback=green!5,
    colframe=green!50,
    breakable,
    title=Données
}

% Métadonnées
\title{Exercices Machine Learning\\
\large Régression Linéaire, KNN, K-means, Arbre de Décision}
\author{AmineGR03}
\date{\today}

\begin{document}

\maketitle

\tableofcontents
\newpage

% ============================================
% PARTIE 1 : RÉGRESSION LINÉAIRE
% ============================================

\section{Régression Linéaire}

\subsection{Exercice 1 : Calcul de la MSE et des paramètres}

\begin{questionbox}
On considère le dataset suivant pour prédire le salaire (Y) en fonction des années d'expérience (X) :

\begin{center}
\begin{tabular}{|c|c|}
\hline
\textbf{Années d'expérience (X)} & \textbf{Salaire (Y)} \\
\hline
1 & 45000 \\
2 & 50000 \\
3 & 60000 \\
4 & 80000 \\
5 & 110000 \\
\hline
\end{tabular}
\end{center}

\textbf{Questions :}
\begin{enumerate}
    \item Calculez manuellement les paramètres $a$ (pente) et $b$ (ordonnée à l'origine) de la régression linéaire $\hat{Y} = aX + b$ en utilisant les formules :
    \begin{align*}
    a &= \frac{n\sum XY - \sum X \sum Y}{n\sum X^2 - (\sum X)^2} \\
    b &= \frac{\sum Y - a\sum X}{n}
    \end{align*}
    
    \item Calculez les prédictions $\hat{Y}$ pour chaque point.
    
    \item Calculez la MSE (Mean Squared Error) :
    $$MSE = \frac{1}{n}\sum_{i=1}^{n}(Y_i - \hat{Y}_i)^2$$
    
    \item Interprétez les valeurs de $a$ et $b$ dans le contexte du problème.
\end{enumerate}
\end{questionbox}

\subsection{Exercice 2 : Gradient Descent}

\begin{questionbox}
On souhaite entraîner un modèle de régression linéaire avec le Gradient Descent.

\textbf{Données initiales :}
\begin{itemize}
    \item $a_0 = 0$ (pente initiale)
    \item $b_0 = 0$ (biais initial)
    \item Taux d'apprentissage $\alpha = 0.01$
    \item Dataset : $(X_1=2, Y_1=4)$, $(X_2=3, Y_2=6)$, $(X_3=4, Y_3=8)$
\end{itemize}

\textbf{Questions :}
\begin{enumerate}
    \item Calculez les dérivées partielles de la fonction de coût MSE par rapport à $a$ et $b$ :
    \begin{align*}
    \frac{\partial MSE}{\partial a} &= \frac{2}{n}\sum_{i=1}^{n}X_i(\hat{Y}_i - Y_i) \\
    \frac{\partial MSE}{\partial b} &= \frac{2}{n}\sum_{i=1}^{n}(\hat{Y}_i - Y_i)
    \end{align*}
    
    \item Effectuez une itération du Gradient Descent :
    \begin{align*}
    a_{new} &= a_{old} - \alpha \frac{\partial MSE}{\partial a} \\
    b_{new} &= b_{old} - \alpha \frac{\partial MSE}{\partial b}
    \end{align*}
    
    \item Calculez la nouvelle valeur de la MSE après cette itération.
    
    \item Expliquez pourquoi le Gradient Descent converge vers les valeurs optimales.
\end{enumerate}
\end{questionbox}

\subsection{Exercice 3 : Questions de cours}

\begin{questionbox}
\textbf{Questions théoriques :}
\begin{enumerate}
    \item Quelle est la différence entre la régression linéaire simple et la régression linéaire multiple ?
    
    \item Expliquez ce qu'est la MSE et pourquoi on l'utilise comme fonction de coût.
    
    \item Qu'est-ce que le Gradient Descent et comment fonctionne-t-il ?
    
    \item Quelle est la différence entre le Gradient Descent batch, stochastique et mini-batch ?
    
    \item Qu'est-ce que le coefficient de détermination $R^2$ et comment l'interpréter ?
    
    \item Expliquez le problème du surapprentissage (overfitting) dans le contexte de la régression linéaire.
\end{enumerate}
\end{questionbox}

\newpage

% ============================================
% PARTIE 2 : KNN (K-Nearest Neighbors)
% ============================================

\section{K-Nearest Neighbors (KNN)}

\subsection{Exercice 1 : Calcul manuel avec petit dataset}

\begin{databox}
On considère un dataset de classification binaire avec 2 caractéristiques (X1, X2) :

\begin{center}
\begin{tabular}{|c|c|c|c|}
\hline
\textbf{Point} & \textbf{X1} & \textbf{X2} & \textbf{Classe} \\
\hline
A & 1 & 2 & Rouge \\
B & 2 & 3 & Rouge \\
C & 3 & 1 & Bleu \\
D & 4 & 2 & Bleu \\
E & 5 & 3 & Bleu \\
\hline
\end{tabular}
\end{center}

On souhaite classifier le point de test $T = (3, 2)$.
\end{databox}

\begin{questionbox}
\textbf{Questions :}
\begin{enumerate}
    \item Calculez la distance euclidienne entre le point T et tous les points du dataset :
    $$d(T, P_i) = \sqrt{(X1_T - X1_i)^2 + (X2_T - X2_i)^2}$$
    
    \item Pour $k=1$, quelle est la classe prédite pour T ? Justifiez.
    
    \item Pour $k=3$, quels sont les 3 plus proches voisins ? Quelle est la classe prédite ?
    
    \item Pour $k=5$, quelle est la classe prédite ? Que remarquez-vous ?
    
    \item Discutez de l'influence du choix de $k$ sur les résultats.
\end{enumerate}
\end{questionbox}

\subsection{Exercice 2 : Matrice de confusion et métriques}

\begin{databox}
On a testé un modèle KNN ($k=3$) sur un dataset de test et obtenu les résultats suivants :

\begin{center}
\begin{tabular}{|c|c|c|}
\hline
\textbf{Point} & \textbf{Classe réelle} & \textbf{Classe prédite} \\
\hline
1 & Positif & Positif \\
2 & Positif & Positif \\
3 & Positif & Négatif \\
4 & Négatif & Positif \\
5 & Négatif & Négatif \\
6 & Négatif & Négatif \\
7 & Négatif & Négatif \\
8 & Positif & Positif \\
9 & Positif & Négatif \\
10 & Négatif & Négatif \\
\hline
\end{tabular}
\end{center}
\end{databox}

\begin{questionbox}
\textbf{Questions :}
\begin{enumerate}
    \item Construisez la matrice de confusion (2x2) avec les valeurs de TP, TN, FP, FN.
    
    \item Expliquez la signification de chaque métrique :
    \begin{itemize}
        \item TP (True Positive)
        \item TN (True Negative)
        \item FP (False Positive)
        \item FN (False Negative)
    \end{itemize}
    
    \item Calculez les métriques suivantes :
    \begin{align*}
    \text{Accuracy} &= \frac{TP + TN}{TP + TN + FP + FN} \\
    \text{Precision} &= \frac{TP}{TP + FP} \\
    \text{Recall (Sensitivity)} &= \frac{TP}{TP + FN} \\
    \text{Specificity} &= \frac{TN}{TN + FP} \\
    \text{F1-Score} &= \frac{2 \times Precision \times Recall}{Precision + Recall}
    \end{align*}
    
    \item Interprétez chaque métrique dans le contexte d'un problème médical (détection de maladie).
    
    \item Dans quel cas privilégieriez-vous la Precision plutôt que le Recall ?
\end{enumerate}
\end{questionbox}

\subsection{Exercice 3 : Questions de cours}

\begin{questionbox}
\textbf{Questions théoriques :}
\begin{enumerate}
    \item Expliquez le principe de l'algorithme KNN.
    
    \item Quels sont les avantages et inconvénients de KNN ?
    
    \item Comment choisir la valeur optimale de $k$ ?
    
    \item Quelle est l'influence de la normalisation des données sur KNN ?
    
    \item Quelle distance utiliser avec KNN ? Comparez la distance euclidienne, manhattan et minkowski.
    
    \item Pourquoi KNN est-il considéré comme un algorithme "lazy learner" ?
    
    \item Expliquez le problème de la malédiction de la dimensionnalité avec KNN.
\end{enumerate}
\end{questionbox}

\newpage

% ============================================
% PARTIE 3 : K-MEANS
% ============================================

\section{K-Means Clustering}

\subsection{Exercice 1 : Calcul manuel de l'algorithme}

\begin{databox}
On considère 6 points dans un espace 2D :

\begin{center}
\begin{tabular}{|c|c|}
\hline
\textbf{Point} & \textbf{Coordonnées (X, Y)} \\
\hline
P1 & (1, 2) \\
P2 & (2, 3) \\
P3 & (3, 1) \\
P4 & (8, 7) \\
P5 & (9, 8) \\
P6 & (7, 9) \\
\hline
\end{tabular}
\end{center}

On souhaite effectuer un clustering avec $k=2$ clusters.
\end{databox}

\begin{questionbox}
\textbf{Questions :}
\begin{enumerate}
    \item Initialisez aléatoirement 2 centroïdes : $C1 = (2, 2)$ et $C2 = (8, 8)$.
    
    \item \textbf{Étape 1 - Assignment :} Assignez chaque point au cluster le plus proche en calculant les distances euclidiennes.
    
    \item \textbf{Étape 2 - Update :} Recalculez les centroïdes comme la moyenne des points de chaque cluster.
    
    \item Répétez les étapes 1 et 2 jusqu'à convergence (les centroïdes ne changent plus).
    
    \item Calculez l'inertie (within-cluster sum of squares) :
    $$Inertie = \sum_{i=1}^{k}\sum_{x \in C_i}||x - \mu_i||^2$$
    où $\mu_i$ est le centroïde du cluster $C_i$.
    
    \item Que se passe-t-il si on choisit $k=3$ au lieu de $k=2$ ?
\end{enumerate}
\end{questionbox}

\subsection{Exercice 2 : Code Python}

\begin{questionbox}
\textbf{Questions de programmation :}
\begin{enumerate}
    \item Écrivez une fonction Python qui calcule la distance euclidienne entre deux points.
    
    \item Écrivez une fonction qui assigne chaque point au cluster le plus proche.
    
    \item Écrivez une fonction qui met à jour les centroïdes.
    
    \item Implémentez l'algorithme K-means complet avec un critère d'arrêt (nombre max d'itérations ou convergence).
    
    \item Utilisez sklearn pour comparer vos résultats avec l'implémentation standard.
\end{enumerate}
\end{questionbox}

\subsection{Exercice 3 : Questions de cours}

\begin{questionbox}
\textbf{Questions théoriques :}
\begin{enumerate}
    \item Expliquez le principe de l'algorithme K-means étape par étape.
    
    \item Quels sont les avantages et inconvénients de K-means ?
    
    \item Comment choisir le nombre optimal de clusters $k$ ? Mentionnez la méthode du coude (elbow method).
    
    \item Pourquoi l'initialisation aléatoire des centroïdes peut-elle affecter les résultats ?
    
    \item Qu'est-ce que l'inertie (within-cluster sum of squares) ?
    
    \item Expliquez le problème des clusters vides dans K-means.
    
    \item Quelle est la différence entre le clustering hiérarchique et K-means ?
    
    \item Pourquoi K-means suppose-t-il que les clusters sont sphériques et de taille similaire ?
\end{enumerate}
\end{questionbox}

\newpage

% ============================================
% PARTIE 4 : ARBRE DE DÉCISION
% ============================================

\section{Arbre de Décision}

\subsection{Exercice 1 : Calcul de l'entropie et du gain d'information}

\begin{databox}
On considère un dataset de classification pour décider si on joue au tennis :

\begin{center}
\begin{tabular}{|c|c|c|c|c|}
\hline
\textbf{Jour} & \textbf{Outlook} & \textbf{Température} & \textbf{Humidité} & \textbf{Jouer} \\
\hline
1 & Ensoleillé & Chaude & Élevée & Non \\
2 & Ensoleillé & Chaude & Élevée & Non \\
3 & Nuageux & Chaude & Élevée & Oui \\
4 & Pluvieux & Douce & Élevée & Oui \\
5 & Pluvieux & Fraîche & Normale & Oui \\
6 & Pluvieux & Fraîche & Normale & Non \\
7 & Nuageux & Fraîche & Normale & Oui \\
8 & Ensoleillé & Douce & Élevée & Non \\
9 & Ensoleillé & Fraîche & Normale & Oui \\
10 & Pluvieux & Douce & Normale & Oui \\
11 & Ensoleillé & Douce & Normale & Oui \\
12 & Nuageux & Douce & Élevée & Oui \\
13 & Nuageux & Chaude & Normale & Oui \\
14 & Pluvieux & Douce & Élevée & Non \\
\hline
\end{tabular}
\end{center}
\end{databox}

\begin{questionbox}
\textbf{Questions :}
\begin{enumerate}
    \item Calculez l'entropie initiale du dataset :
    $$Entropie(S) = -\sum_{i=1}^{c}p_i \log_2(p_i)$$
    où $p_i$ est la proportion de la classe $i$ dans l'ensemble $S$.
    
    \item Calculez l'entropie pour chaque attribut (Outlook, Température, Humidité) :
    $$Entropie(S, A) = \sum_{v \in Values(A)}\frac{|S_v|}{|S|}Entropie(S_v)$$
    où $S_v$ est le sous-ensemble de $S$ où l'attribut $A$ a la valeur $v$.
    
    \item Calculez le gain d'information pour chaque attribut :
    $$Gain(S, A) = Entropie(S) - Entropie(S, A)$$
    
    \item Quel attribut doit être choisi comme racine de l'arbre ? Justifiez.
    
    \item Construisez l'arbre de décision complet en répétant le processus pour chaque nœud.
    
    \item Utilisez l'arbre pour classifier un nouveau jour : (Ensoleillé, Chaude, Normale).
\end{enumerate}
\end{questionbox}

\subsection{Exercice 2 : Calcul du Gini Impurity}

\begin{questionbox}
\textbf{Questions :}
\begin{enumerate}
    \item Calculez l'impureté de Gini initiale du dataset de l'exercice 1 :
    $$Gini(S) = 1 - \sum_{i=1}^{c}p_i^2$$
    
    \item Calculez le Gini Gain pour chaque attribut :
    $$Gini(S, A) = \sum_{v \in Values(A)}\frac{|S_v|}{|S|}Gini(S_v)$$
    $$GiniGain(S, A) = Gini(S) - Gini(S, A)$$
    
    \item Comparez les résultats obtenus avec l'entropie et le Gini. Quel critère choisiriez-vous et pourquoi ?
\end{enumerate}
\end{questionbox}

\subsection{Exercice 3 : Questions de cours}

\begin{questionbox}
\textbf{Questions théoriques :}
\begin{enumerate}
    \item Expliquez le principe de construction d'un arbre de décision.
    
    \item Quelle est la différence entre l'entropie et l'impureté de Gini ?
    
    \item Qu'est-ce que le gain d'information et comment l'utilise-t-on ?
    
    \item Expliquez le problème du surapprentissage (overfitting) avec les arbres de décision.
    
    \item Quelles sont les techniques d'élagage (pruning) pour éviter le surapprentissage ?
    
    \item Quelle est la différence entre le pré-élagage (pre-pruning) et le post-élagage (post-pruning) ?
    
    \item Comment gérer les valeurs manquantes dans un arbre de décision ?
    
    \item Quels sont les avantages et inconvénients des arbres de décision ?
    
    \item Expliquez le concept de forêt aléatoire (Random Forest) et son lien avec les arbres de décision.
\end{enumerate}
\end{questionbox}

\newpage

% ============================================
% CONCLUSION
% ============================================

\section{Conclusion}

Ces exercices couvrent les concepts fondamentaux du Machine Learning :
\begin{itemize}
    \item \textbf{Régression Linéaire} : Calcul des paramètres, MSE, Gradient Descent
    \item \textbf{KNN} : Classification avec calculs manuels, matrice de confusion, métriques
    \item \textbf{K-means} : Clustering avec calculs itératifs et implémentation
    \item \textbf{Arbre de Décision} : Entropie, gain d'information, construction d'arbre
\end{itemize}

\textbf{Bon courage pour vos révisions !}

\end{document}

